%%%%%%%%%%%%%%%%%
% This is an sample CV template created using altacv.cls
% (v1.6.5, 3 Nov 2022) written by LianTze Lim (liantze@gmail.com). Compiles with pdfLaTeX, XeLaTeX and LuaLaTeX.
%
%% It may be distributed and/or modified under the
%% conditions of the LaTeX Project Public License, either version 1.3
%% of this license or (at your option) any later version.
%% The latest version of this license is in
%%    http://www.latex-project.org/lppl.txt
%% and version 1.3 or later is part of all distributions of LaTeX
%% version 2003/12/01 or later.
%%%%%%%%%%%%%%%%

%% Use the "normalphoto" option if you want a normal photo instead of cropped to a circle
% \documentclass[10pt,a4paper,normalphoto]{altacv}

\documentclass[10pt,a4paper,ragged2e,withhyper]{altacv}
%% AltaCV uses the fontawesome5 and packages.
%% See http://texdoc.net/pkg/fontawesome5 for full list of symbols.

% Change the page layout if you need to
\geometry{left=1cm,right=1cm,top=0.8cm,bottom=0.8cm,columnsep=1cm}

% The paracol package lets you typeset columns of text in parallel
\usepackage{paracol}

% Change the font if you want to, depending on whether
% you're using pdflatex or xelatex/lualatex
\ifxetexorluatex
  % If using xelatex or lualatex:
  \setmainfont{Roboto Slab}
  \setsansfont{Lato}
  \renewcommand{\familydefault}{\sfdefault}
\else
  % If using pdflatex:
  \usepackage[rm]{roboto}
  \usepackage[defaultsans]{lato}
  % \usepackage{sourcesanspro}
  \renewcommand{\familydefault}{\sfdefault}
\fi

% Change the colours if you want to
\definecolor{SlateGrey}{HTML}{2E2E2E}
\definecolor{LightGrey}{HTML}{666666}
\definecolor{DarkPastelRed}{HTML}{450808}
\definecolor{PastelRed}{HTML}{8F0D0D}
\definecolor{GoldenEarth}{HTML}{E7D192}
\colorlet{name}{black}
\colorlet{tagline}{PastelRed}
\colorlet{heading}{DarkPastelRed}
\colorlet{headingrule}{GoldenEarth}
\colorlet{subheading}{PastelRed}
\colorlet{accent}{PastelRed}
\colorlet{emphasis}{SlateGrey}
\colorlet{body}{LightGrey}

% Change some fonts, if necessary
\renewcommand{\namefont}{\Huge\rmfamily\bfseries}
\renewcommand{\personalinfofont}{\footnotesize}
\renewcommand{\cvsectionfont}{\LARGE\rmfamily\bfseries}
\renewcommand{\cvsubsectionfont}{\large\bfseries}


% Change the bullets for itemize and rating marker
% for \cvskill if you want to
\renewcommand{\itemmarker}{{\small\textbullet}}
\renewcommand{\ratingmarker}{\faCircle}

%% Bibliography was configured here previously, but is unused.

\begin{document}
\name{José de Deus e Silva Neto}
\tagline{Recherche du premier emploi — à partir de septembre 2027}
%% You can add multiple photos on the left or right
% \photoR{2.8cm}{Globe_High}
% \photoL{2.5cm}{Yacht_High,Suitcase_High}

\personalinfo{%
  % Not all of these are required!
  \email{nectxx@gmail.com}
  \phone{+55 (91) 98565-6002}
  % \mailaddress{21, Avenue de Verdun}
  \location{Belém, Pará, Brasil}
  %\homepage{www.homepage.com}
  %\twitter{@twitterhandle}
  \linkedin{jose-de-deus-neto}
  \github{josedsneto}
  %\orcid{0000-0000-0000-0000}
  %% You can add your own arbitrary detail with
  %% \printinfo{symbol}{detail}[optional hyperlink prefix]
  % \printinfo{\faPaw}{Hey ho!}[https://example.com/]
  %% Or you can declare your own field with
  %% \NewInfoFiled{fieldname}{symbol}[optional hyperlink prefix] and use it:
  % \NewInfoField{gitlab}{\faGitlab}[https://gitlab.com/]
  % \gitlab{your_id}
  %%
  %% For services and platforms like Mastodon where there isn't a
  %% straightforward relation between the user ID/nickname and the hyperlink,
  %% you can use \printinfo directly e.g.
  % \printinfo{\faMastodon}{@username@instace}[https://instance.url/@username]
  %% But if you absolutely want to create new dedicated info fields for
  %% such platforms, then use \NewInfoField* with a star:
  % \NewInfoField*{mastodon}{\faMastodon}
  %% then you can use \mastodon, with TWO arguments where the 2nd argument is
  %% the full hyperlink.
  % \mastodon{@username@instance}{https://instance.url/@username}
}

\makecvheader
%% Depending on your tastes, you may want to make fonts of itemize environments slightly smaller
\AtBeginEnvironment{itemize}{\small\setlength{\itemsep}{0pt}\setlength{\parskip}{0pt}\setlength{\parsep}{0pt}}

%% Set the left/right column width ratio to 6:4.
\columnratio{0.6}

% Start a 2-column paracol. Both the left and right columns will automatically
% break across pages if things get too long.
\begin{paracol}{2}
\cvsection{Expérience}

\cvevent{Développeur Logiciel}{Pulsar Audio SARL}{Février 2026 --- Août 2026}{Grenoble --- France}

\begin{itemize}
\item Conception d’un gestionnaire de plugins audio à l’aide de Rust et du framework Tauri.
\item Adoption d’une méthodologie de développement accéléré («~Vibe Coding~»).
\item Validation, tests et documentation du logiciel.
\end{itemize}

\divider

\cvevent{Ingénieur en Cybersécurité}{Grenoble INP - ESISAR, UGA | LCIS}{Septembre 2025 --- Janvier 2026}{Valence --- France}

\begin{itemize}
\item Amélioration d’une plateforme d’authentification matérielle basée sur des PUF (Physical Unclonable Function).
\item Vérification de la robustesse face aux variations de température et analyse de la précision d’identification.
\item Programmation de prototypes à l’aide de la chaîne d’outils ChipWhisperer.
\end{itemize}

\divider

\cvevent{Développeur Firmware}{Grenoble INP - ESISAR, UGA | Olyntec}{Janvier 2025 --- Juin 2025}{Valence --- France}

\begin{itemize}
\item Rédaction de règles comportementales en C pour optimiser l’autonomie énergétique via la sélection de réseaux.
\item Gestion des trames sur les réseaux LTE-M et LoRaWAN.
\item Analyse des modes de basse consommation eDRX et PSM pour un prototype IoT.
\item Géolocalisation à l’aide de la bibliothèque Uber/H3.
\end{itemize}

\divider

\cvevent{Chercheur 5G}{LASSE/UFPA | Ericsson}{Juin 2022 --- Septembre 2024}{Belém --- Brésil}

\begin{itemize}
\item Contribution à une version interne du projet 
\printinfoo{\faGithub}{PTP-DAL}[https://github.com/lasseufpa/ptp-dal], une bibliothèque Python d’analyse de jeux de données IEEE 1588 (Precision Time Protocol).
\item Étude de la synchronisation d'un Fronthaul 5G utilisant des schémas TDD.
%\item Deep neural networks programming with PyTorch.
\end{itemize}

% \cvevent{Scholarship Student}{PET --- UFPA}{2021 --- 2022}{Belém --- Brazil}

% \begin{itemize}
% \item Lectures on Matlab, Arduino and AutoCad for electrical engineering students.
% \item Writing of academic works for national meetings and conferences.

% \end{itemize}

% \cvevent{D\'eveloppeur PCB}{Robotech/UFPA}{Février 2020 --- Février 2021}{Belém --- Brésil}

% \begin{itemize}
% % \item KiCad and PCB design for competitive robotics in mini-sumo category.
% \item Conception de circuits électroniques pour robots mini-sumo à l’aide de l’outil KiCad.
% % \item Power management techniques.
% \end{itemize}

\cvsection{Recherche}

\cvevent{Synchronization Algorithms for 5G and 6G Smarthauling}{CNPq/UFPA}{}{}
\begin{itemize}
  \item Étude de techniques de synchronisation OTA pour les réseaux de nouvelle génération.
  \item Application du Deep Learning à l’amélioration de la synchronisation temporelle dans les architectures C-RAN.
  \item Distribution d’horloge à l’aide du projet Linux PTP.
\end{itemize}

\switchcolumn

\cvsection{Formation}
\cvevent{Double Diplôme en Électronique, Informatique et Systèmes}{Grenoble INP - ESISAR, UGA}
{Septembre 2024 -- Août 2026}{}

\divider

\cvevent{Diplôme d'Ingénieur en Génie Électrique}{Université Fédérale du Pará (UFPA)}{Février 2020 -- Présent}{}

% \newpage

% \divider

% \cvevent{PROEX - Pro-Rectory of Extension}{PROEX - UFPA}{2023}{}
% \begin{itemize}
%     \item \href{https://ufpabr-my.sharepoint.com/:b:/r/personal/jose_silva_neto_itec_ufpa_br/Documents/Documentos/Certificados/certificado_5g_ufpa_mar%C3%A7o_Jos%C3%A9%20de%20Deus%20e%20Silva%20Neto.pdf?csf=1&web=1&e=a5ibI0}{Training in 5G Technologies - Huawei adapted content HCIA - 5G}.
% \end{itemize}

% \cvevent{PROEX - Pro-Rectory of Extension}{PROEX - UFPA}{2023}{}
% \begin{itemize}
%     \item Taken Training in 5G Technologies - Huawei adapted content
% HCIA - 5G.
% \end{itemize}
% \divider
% \divider

% \cvevent{Project 2}{Funding agency/institution}{Project duration}{}
% A short abstract would also work.

% \medskip


% use ONLY \newpage if you want to force a page break for
% ONLY the current column
% \newpage

% \cvsection{Publications}

% %% Specify your last name(s) and first name(s) as given in the .bib to automatically bold your own name in the publications list. 
% %% One caveat: You need to write \bibnamedelima where there's a space in your name for this to work properly; or write \bibnamedelimi if you use initials in the .bib
% %% You can specify multiple names, especially if you have changed your name or if you need to highlight multiple authors. 
% \mynames{Lim/Lian\bibnamedelima Tze,
%   Wong/Lian\bibnamedelima Tze,
%   Lim/Tracy,
%   Lim/L.\bibnamedelimi T.}
% %% MAKE SURE THERE IS NO SPACE AFTER THE FINAL NAME IN YOUR \mynames LIST

% \nocite{*}

% \printbibliography[heading=pubtype,title={\printinfo{\faBook}{Books}},type=book]

% \divider

% \printbibliography[heading=pubtype,title={\printinfo{\faFile*[regular]}{Journal Articles}},type=article]

% \divider

% \printbibliography[heading=pubtype,title={\printinfo{\faUsers}{Conference Proceedings}},type=inproceedings]

%% Switch to the right column. This will now automatically move to the second
%% page if the content is too long.
% \switchcolumn

\cvsection{Profil \& Objectifs}

% \begin{quote}

Étudiant passionné par la programmation, l’intelligence artificielle, les systèmes embarqués et l’électronique. 
Au cours de mes deux années d’échange en France, où je prépare un double diplôme, j’ai participé à 
divers projets et réalisé un stage pratique en développement logiciel. Auparavant, j’ai acquis 
deux ans d’expérience au sein d’un laboratoire de R\&D. Ce parcours m’a apporté une forte 
capacité d’adaptation, une grande aisance pour le travail en équipe ainsi qu’une réelle 
autonomie, en plus de bonnes compétences rédactionnelles en anglais. Motivé et curieux, je 
cherche constamment à apprendre et à relever de nouveaux défis.

%I'm curious and very excited about the new technologies in the AI era. As an international student a could accumulate
%a very rich culture. Along my career I worked on software development, embedded  systems, IoT and ML. That includes 2 years
%of experience in R&D laboratory and 6 months being followed by a company in a industrial project. That allowed me to work
%alone or as a team applying agile methods. Personally, I like being challenged and participating in multicultural %environments, where I can learn about the world and myself.

%I'm highly interested in fields related to software development, embedded systems, telecommunication and artificial
%intelligence. Beyond my professional skills, I consider empathy and communication strong points of my personality. I'm
%always open to new challenges and opportunities and I am flexible to work as a team or by myself.
% \end{quote}

% \cvsection{Most Proud of}

% \cvachievement{\faTrophy}{Fantastic Achievement}{and some details about it}

% \divider

% \cvachievement{\faHeartbeat}{Another achievement}{more details about it of course}

% \divider

% \cvachievement{\faHeartbeat}{Another achievement}{more details about it of course}

\cvsection{Compétences}

\cvtag{Rust}
\cvtag{TypeScript}
\cvtag{Python}
\cvtag{Linux} \\
\cvtag{C/C++} 
\cvtag{Git} 
%\cvtag{PyTorch}\\
\cvtag{VHDL}
\cvtag{MATLAB} \\
% \cvtag{Arduino} \\
\cvtag{Java}
\cvtag{Docker}
\cvtag{DSP} 
\cvtag{KiCad}
% \cvtag{Lua}
% \cvtag{JavaScript}\\
% \cvtag{IEEE 1588}
% \cvtag{Fronthaul 5G} 
% \cvtag{FPGA}
% \cvtag{Vivado} \\
% \cvtag{C/C++}\\

\cvsection{Langues}

\cvskill{Portugais - Natif}{5} %% Supports X.5 values.
\divider

\cvskill{Anglais - TOEIC B2}{4}
\divider

\cvskill{Français - TCF B2}{4}
\divider

%\cvskill{Allemand}{1.5} %% Supports X.5 values.

% \cvsection{Publications}
% \begin{itemize}
%     \item Neto, José. Silva, André. Campos, Matheus. Matos, Wendler. Silva, Orlando. (2021). Machine Learning models applied to breast cancer detection. \href{https://sbic.org.br/eventos/cbic_2021/cbic2021-154/}{10.21528/CBIC2021-154.} 
% \end{itemize}

%% Yeah I didn't spend too much time making all the
%% spacing consistent... sorry. Use \smallskip, \medskip,
%% \bigskip, \vspace etc to make adjustments.
%\medskip

% \divider

\cvsection{Références}

% \cvref{name}{email}{mailing address}
\cvref{Prof.\ Aldebaro Klautau}{\href{https://www.lasse.ufpa.br/en}{LASSE/UFPA}}{aldebaro@ufpa.br}
{66075110 - Belém, PA - Brésil}

\end{paracol}


\end{document}
